\subsection{开集的外测度}

给定开集 \(G \subset X\) ，其特征函数 \(\varphi_{G}\) 在 X 上是下半连续的（GT, IV, §6, No.~2, 命题 1 的推论）。因此可以作下述定义：

定义 2. --- 给定 X 上的正测度 \(\mu\) ，对每个开集 \(G \subset X\) ，上积分 \(\mu^{*}(\varphi_{G})\) 称为 G 的外测度，记作 \(\mu^{*}(G)\) 。

这样，开集 G 的外测度是一个 \(\geqslant0\) 的数，有限或等于 \(+\infty\) 。显然 \(\mu^{*}(\varnothing)=0\) 。此外，\(\mu^{*}(X)=\|\mu\|\) ，这由 Ch. III, §1, No.~8 的公式 (22) 得出。

命题 5. --- 相对紧的开集 G 的外测度是有限的。

事实上，此时存在函数 \(f \in \mathcal{K}_+\) 使 \(\varphi_{\mathrm{G}} \leqslant f\) （Ch. III, §1, No.~2, 引理 1），由此

\[
\mu^ {*} (\mathrm{G}) = \mu^ {*} (\varphi_ {\mathrm{G}}) \leqslant \mu^ {*} (f) = \mu (f) <   + \infty .
\]

外测度有限的开集并不总是相对紧的（习题 3）。

命题 6. --- 若 \(G_1\) 与 \(G_2\) 是满足 \(G_1 \subset G_2\) 的两个开集，则 \(\mu^*(G_1) \leqslant \mu^*(G_2)\) 。

事实上，关系 \(\mathbf{G}_1\subset \mathbf{G}_2\) 等价于 \(\varphi_{\mathrm{G_1}}\leqslant \varphi_{\mathrm{G_2}}\)

命题 7. --- 设 G 是 X 的开子集所成的一个集，它按关系 ⊂ 有向；则

\[
\mu^ {*} \left(\bigcup_ {G \in \mathfrak {G}} G\right) = \sup _ {G \in \mathfrak {G}} \mu^ {*} (G).\tag{5}
\]

诸函数 \(\varphi_{\mathbf{G}}\) 构成一个（按 \(\leqslant\) ）有向集，且都属于 \(\mathcal{I}_{+}\) ；它们的上包络是诸集 \(\mathbf{G} \in \mathfrak{G}\) 之并的特征函数；因此本命题是定理 1 的推论。

命题 8. --- 设 \((G_{\iota})_{\iota \in I}\) 是任意一族开集；则

\[
\mu^ {*} \left(\bigcup_ {\iota \in I} G _ {\iota}\right) \leqslant \sum_ {\iota \in I} \mu^ {*} (G _ {\iota}).\tag{6}
\]

此外，若诸 \(G_{\iota}\) 两两不相交，则

\[
\mu^ {*} \left(\bigcup_ {\iota \in I} G _ {\iota}\right) = \sum_ {\iota \in I} \mu^ {*} (G _ {\iota}).\tag{7}
\]

事实上，若 \(G = \bigcup_{\iota \in I} G_{\iota}\) ，则 \(\varphi_{G} = \sup_{\iota \in I} \varphi_{G_{\iota}} \leqslant \sum_{\iota \in I} \varphi_{G_{\iota}}\) ；当诸 \(G_{\iota}\) 两两不相交时，\(\varphi_{G} = \sum_{\iota \in I} \varphi_{G_{\iota}}\) ；因此本命题是命题 2 与命题 3 的推论。

例子. --- 设 X = R ，并设 \(\mu\) 是 R 上的 Lebesgue 测度（Ch. III, §1, No.~3）；我们来确定开区间 G = {]}a, b{[} （ \(-\infty \leqslant a < b \leqslant +\infty\) ）的外测度。首先假设 a 与 b 都有限。对 \(K_{+}\) 中每个满足 \(f \leqslant \varphi_{G}\) 的函数 f ，由平均定理（FRV, II, §1, No.~5, 命题 6）有

\[
\int_ {- \infty} ^ {+ \infty} f (x) d x = \int_ {a} ^ {b} f (x) d x \leqslant b - a,
\]

由此 \(\mu^{*}(\mathrm{G})\leqslant b - a\) 。另一方面，对每个 \(\varepsilon >0\) ，存在函数 \(f\in \mathcal{K}_{+}\) 满足 \(f\leqslant \varphi_{\mathrm{G}}\) ，且 \(f(x) = 1\) 对 \(a + \varepsilon \leqslant x\leqslant b - \varepsilon\) 成立，由此 \(\mu^{*}(G) \geqslant b - a - 2\varepsilon\) ；由于 \(\varepsilon\) 是任意的，故 \(\mu^{*}(G) = b - a\) ；换言之，G 的外测度等于它的长度。这一结果立即推广到 G 为无界开区间的情形，因为此时 G 含有长度任意大的有界开区间；从而在此情形 \(\mu^{*}(G) = +\infty\) 。

现在设 G 是 R 中任一开集；G 是可数个（有限个或无穷个）两两不相交的开区间 \({]}a_{k}, b_{k}{[}\) 的并（GT, IV, §2, No.~5, 命题 2），因此

\[
\mu^ {*} (\mathrm{G}) = \sum_ {k} (b _ {k} - a _ {k})
\]

（命题 8）；换言之：

命题 9. --- 对 R 上的 Lebesgue 测度而言，R 中开集的外测度等于其诸连通分支的长度之和。

特别地，注意若 G 是 R 中满足 \(\mu^{*}(G)=0\) 的开集，则 G 是空集。

